\chapter{Morphismes lisses: propri\'et\'es~de~prolongement}
\label{III}

\section{Homomorphismes formellement lisses}
\label{III.1}

Dans~\ref{II}, nous nous sommes born\'es \`a la consid\'eration
d'homomorphismes de type fini, et par cons\'equent, dans les
homomorphismes locaux~$A\to B$ d'anneaux locaux, au cas o\`u~$B$ est
isomorphe \`a une alg\`ebre localis\'ee d'une $A$-alg\`ebre de
type fini. Ce cas est insuffisant pour diverses applications,
notamment en g\'eom\'etrie formelle ou en g\'eom\'etrie
analytique. Par exemple, l'anneau de s\'eries
formelles~$B=A[[t_1,\dots,t_n]]$ a (du point de vue de la
g\'eom\'etrie formelle) les propri\'et\'es d'une alg\`ebre
lisse sur~$A$. En g\'eom\'etrie analytique, il en est de m\^eme
de l'anneau local d'un point~$(y,z)$ d'un produit $Y\times\CC^n$
consid\'er\'e comme alg\`ebre sur l'anneau local de~$y$;
d'ailleurs, la compl\'et\'ee de cette alg\`ebre est isomorphe
\`a l'alg\`ebre des s\'eries formelles en~$n$
ind\'etermin\'ees sur le compl\'et\'e de l'anneau de
base~$\cal{O}_x$. C'est ce qui conduit \`a poser la d\'efinition
qui suit.

\begin{definition}
\label{III.1.1}
Soit~$u\colon A\to B$ un homomorphisme local d'anneaux locaux
(noeth\'eriens, on le rappelle). On 
suppose~$\kres(B)$ fini sur~$\kres(A)$.
On dit que~$u$ est un \emph{homomorphisme
formellement lisse},
\index{formellement lisse (homomorphisme, alg\`ebre)|hyperpage}\index{lisse (formellement)|hyperpage}ou que l'alg\`ebre~$B$ est \emph{formellement lisse sur}~$A$, s'il
existe une $\overline{A}$-alg\`ebre locale
finie~$A'$,
libre sur~$\overline{A}$, telle que les composants locaux de l'anneau
semi-local~$\overline{B}\otimes_{\overline{A}} A'=B'$ soient
$A'$-isomorphes \`a des alg\`ebres des s\'eries formelles
sur~$A'$\footnote{Pour une d\'efinition plus g\'en\'erale et
plus conceptuelle, motiv\'ee par~\ref{III.2.1} ci-dessous,
cf.\ EGA~$\mathrm{0}_{\mathrm{IV}}$~19.3.1.}.
\end{definition}
(On d\'enote par~$\overline{A}$, $\overline{B}$ les anneaux
compl\'et\'es de~$A$, $B$). Comme~$B'$ est fini et libre
sur~$\overline{B}$, c'est bien un anneau semi-local, compos\'e
direct d'anneaux locaux complets, dont chacun est encore un module
libre sur~$\overline{B}$, donc a m\^eme dimension que~$\overline{B}$
donc que~$B$. Il s'ensuit que le nombre de variables~$t_i$ dans les
anneaux de s\'eries formelles envisag\'es dans~\ref{III.1.1} est
\'egal \`a $\dim \overline{B}-\dim \overline{A}=\dim B-\dim A$, et
en particulier ind\'ependant du composant local
choisi. On voit tout de suite que c'est aussi la dimension de l'anneau
$B\otimes k=B/\goth{m} B$, o\`u~$k=A/\goth{m}$ est le corps
r\'esiduel de~$A$; on l'appellera la \emph{dimension relative
de}~$B$
\index{dimension relative|hyperpage}\emph{par rapport à}~$A$.

\begin{remarques}
\label{III.1.2}
Il est \'evident que la d\'efinition~\ref{III.1.1} ne d\'epend
que de l'homomorphisme sur les compl\'et\'es $\overline{A} \to
\overline{B}$ d\'eduit de~$A\to B$, ce qui justifie dans une
certaine mesure la terminologie. Nous nous repentons ici de la
d\'efinition~I~\ref{I.3.2}~b) et I~\ref{I.4.1}~b), qui risque
d'induire en erreur, et pr\'ef\'erons dire \og formellement non
ramifi\'e\fg
\index{formellement non ramifi\'e (\resp net)|hyperpage}et \og formellement \'etale\fg
\index{formellement etale@formellement \'etale|hyperpage}dans les cas envisag\'es dans ces d\'efinitions, r\'eservant la
terminologie \og non ramifi\'ee\fg et \og \'etale\fg au cas o\`u $B$
est localis\'ee d'une $A$-alg\`ebre de type fini\footnote{Ou
mieux, \og essentiellement non ramifi\'e\fg \resp \og essentiellement
\'etale\fg, comparer EGA IV 18.6.1.}.
\index{essentiellement etale@essentiellement \'etale|hyperpage}\index{essentiellement non ramifi\'e (\resp net)|hyperpage}Le lecteur v\'erifiera aussit\^ot que \og formellement \'etale\fg
\'equivaut \`a \og formellement lisse et quasi-fini\fg. Enfin,
signalons qu'il y a une d\'efinition raisonnable de \og formellement
lisse\fg sans aucune hypoth\`ese pr\'ealable sur l'extension
r\'esiduelle
$\kres(B)/\kres(A)$
(suppos\'ee ici finie),
englobant entre autres les homomorphismes locaux $A\to B$ tels que $B$
soit \emph{plat} sur~$A$ et $B/\goth{m}B$ une \emph{extension
séparable} de~$A/\goth{m} =k$ (pas n\'ecessairement de type
fini); par exemple, un $p$-anneau de Cohen est formellement
lisse
sur
l'anneau des entiers $p$-adiques. C'est la propri\'et\'e de
rel\`evement des homomorphismes (comparer~\ref{III.2.1}) qui doit
devenir d\'efinition dans ce cas g\'en\'eral. Pour les
applications que nous avons en vue, le cas trait\'e dans la
d\'efinition~\ref{III.1.1} nous suffira; par la suite, dans
\og formellement lisse\fg nous sous-entendrons \og \`a extension
r\'esiduelle finie\fg.
\end{remarques}

\begin{lemme}
\label{III.1.3}
Si $B$ est formellement lisse sur~$A$, $B$ est plat sur~$A$.
\end{lemme}


Comme la platitude est invariante par compl\'etion, on peut supposer
$A$ et $B$ complets. Comme la platitude est invariante par extension
plate locale (donc fid\`element plate) de l'anneau de base, on est
ramen\'e en vertu de d\'efinition~\ref{III.1.1} au cas o\`u $B$
est une alg\`ebre de s\'eries formelles sur~$A$. Mais alors en
tant que $A$-module, $B$ est isomorphe \`a un produit de~$A$-modules
isomorphes \`a $A$, donc (l'anneau de base $A$ \'etant
noeth\'erien) est $A$-plat comme produit de~$A$-modules plats.


\medskip
Mettons-nous sous les conditions de~\ref{III.1.1}. Comme les
extensions
r\'esiduelles des composants locaux de~$B'$ sur~$A'$ sont triviales,
il s'ensuit que $L\otimes_k k'$ est une $k'$-alg\`ebre artinienne
dont les composants locaux ont des extensions r\'esiduelles
triviales (o\`u $L,\, k,\, k'$ sont les corps r\'esiduels de~$A,\,
B,\, A'$). Cette condition n\'ecessaire pour que l'extension finie
libre $A'$ satisfasse la condition \'enonc\'ee dans~\ref{III.1.1}
est aussi suffisante, comme il r\'esulte aussit\^ot de
\ref{III.1.4}~(i) et~\ref{III.1.5} ci-dessous.

\begin{proposition}
\label{III.1.4}
Soit $A\to B$ un homomorphisme local d'anneaux locaux, \`a extension
r\'esiduelle finie
et
soit $A'$ une $A$-alg\`ebre finie locale sur~$A$, de sorte que
$B'=B\otimes_A A'$ est finie sur~$B$, donc un anneau semi-local
\'egalement noeth\'erien. \textup{(i)} Si $B$ est formellement lisse
sur~$A$, alors les localis\'es de B en ses id\'eaux maximaux sont
formellement lisses sur~$A'$. \textup{(ii)} La r\'eciproque est vraie
si $A'$ est libre sur~$A$.
\end{proposition}

On est aussit\^ot ramen\'e au cas o\`u $A$, $B$ sont complets.

(i) Soit $A''$ une extension finie libre locale de~$A$ telle que les
composants locaux de~$B''=B\otimes_A A''$ soient des alg\`ebres de
s\'eries formelles sur~$A''$. Faisant l'extension des scalaires
$A''\to A''\otimes_A A' \to A'''$, o\`u $A'''$ est un composant
local de~$A''\otimes_A A'$, on voit que les composants locaux de
$B''\otimes_{A''} A'''=B\otimes_A A'''$ sont des alg\`ebres de
s\'eries formelles sur~$A'''$. Or on a aussi
$B\otimes_AA'''=(B\otimes_A A')\otimes_{A'} A'''=B'\otimes_{A'} A'''$,
d'autre
part comme $A''$ est libre sur~$A$, $A''\otimes_A A'$ est libre sur~$A'$ et il en est de m\^eme par suite de~$A'''$ qui en est facteur
direct, ce qui prouve que $B'$ est formellement lisse sur~$A'$.

(ii) Soit $A''$ une $A'$-alg\`ebre finie libre locale telle que les
composants locaux de~$B'\otimes_{A'} A'' = B\otimes_A A''$ soient des
alg\`ebres de s\'eries formelles sur~$A''$. Comme $A'$ est libre
sur~$A$, $A''$ l'est aussi, donc $B$ est formellement \emph{lisse}
sur~$A$.

\begin{proposition}
\label{III.1.5}
Soit $A\to B$ un homomorphisme local d'anneaux locaux, \`a extension
r\'esiduelle \emph{triviale}. Pour que $B$ soit formellement lisse
sur~$A$, il faut et il suffit que $\overline{B}$ soit isomorphe \`a
une alg\`ebre de s\'eries formelles sur~$\overline{A}$.
\end{proposition}

Il n'y a \`a prouver que la n\'ecessit\'e, et on peut supposer
$A$ et $B$ complets. Soit $\goth{m}$ ($\goth{n}$) l'id\'eal maximal
de~$A$ ($B$) et soient $t_1,\dots, t_n$ des \'el\'ements de
$\goth{n}$ qui d\'efinissent une base de l'espace vectoriel
$$
(\goth{n}/\goth{n}^2)/ \Im (\goth{m}/\goth{m}^2)=
\goth{n}/(\goth{n}^2 +\goth{m}B)
$$
Ces \'el\'ements d\'efinissent donc un homomorphisme de
$A$-alg\`ebres locales
$$
B_1= A[[t_1,\dots,t_n]] \to B
$$
prouvons
que c'est un isomorphisme. Il suffit de prouver que pour toute
puissance $\goth{m}^q$ de~$\goth{m}$, on obtient un isomorphisme en
r\'eduisant mod~$\goth{m}^q$ (puisque $B_1$ et $B$ sont les
limites projectives des anneaux correspondants r\'eduits
mod~$\goth{m}^q$, $q$ variable). Comme $B$ et $B_1$ sont des
$A$-modules plats donc les gradu\'es associ\'es \`a la
filtration $\goth{m}$-adique s'obtiennent en tensorisant par $\gr
(A)$, sur~$k=A/\goth{m}$, les anneaux $B_1/\goth{m}B_1$
\resp $B/\goth{m}B$, on est ramen\'e \`a montrer que
$B_1/\goth{m}B_1\to B/\goth{m}B$ est un isomorphisme. Compte tenu
de~\ref{III.1.3}, on est ainsi ramen\'e au cas o\`u $A$ est un
\emph{corps} $k$. D'autre part, si $A'$ est $A$-alg\`ebre finie
libre locale telle que $B\otimes_A A'$ soit une alg\`ebre de
s\'eries formelles sur~$A'$ (N.B. cette alg\`ebre est locale,
puisque l'extension r\'esiduelle de~$B$ sur~$A$ est triviale), pour
prouver que $B_1\to B$ est un isomorphisme, il suffit de prouver que
$B_1\otimes_A A'\to B\otimes_A A'$ l'est. Cela nous ram\`ene au cas
o\`u $B$ est d\'ej\`a une alg\`ebre de s\'eries formelles
(il fallait commencer par cette r\'eduction, avant de se ramener au
cas d'un corps de base). Mais alors $B$ est un anneau local
r\'egulier \`a corps de repr\'esentants $k$, et il est bien
connu (et imm\'ediat par consid\'eration des gradu\'es
associ\'es \`a la filtration $\goth{n}_1 $-adique et
$\goth{n}$-adique sur~$B_1$ et $B$) que $B_1 \to B$ est un
isomorphisme, ce qui ach\`eve la d\'emonstration.

\begin{corollaire}
\label{III.1.6}
Si $B$ est formellement lisse sur~$A$, alors il existe une
$A$-alg\`ebre finie locale $A'$ telle que les composants locaux de
$\overline{B}\otimes_{\overline{A}}
\overline{A'}=\overline{(B\otimes_A A')}$ soient isomorphes \`a des
alg\`ebres de s\'eries formelles sur~$\overline{A'}$.
\end{corollaire}
En effet, si $L/k$ est l'extension r\'esiduelle de~$B/A$, on
consid\`ere une extension $k'/k$, telle que les extensions
r\'esiduelles dans la $k'$-alg\`ebre $L\otimes_k k'$ soient
triviales. Soit $A'$ une alg\`ebre finie libre sur~$A$ telle que
$A'/\goth{m}A'=k'$ (on sait qu'il en existe, par exemple en se
ramenant de proche en proche au cas o\`u $k'/k$ est monog\`ene, et
alors on rel\`eve dans $A$ les coefficients du polyn\^ome minimal
d'un g\'en\'erateur de~$k'$ sur~$k$). Elle est locale. Alors
$B\otimes_A A'$ a en ses id\'eaux maximaux des extensions
r\'esiduelles triviales au-dessus de celle $k'$ de~$A'$, et on
ach\`eve \`a l'aide de~\ref{III.1.5}.

\begin{corollaire}
\label{III.1.7}
Soit $A\to B$ un homomorphisme local d'anneaux locaux. Pour que $B$
soit formellement lisse sur~$A$, il faut et il suffit que $B$ soit
plat sur~$A$ et que $B/\goth{m}B$ soit formellement lisse sur~$k=A/\goth{m}$.
\end{corollaire}
Faisant une extension finie libre locale $A'$ convenable de~$A$ et
utilisant \ref{III.1.4}~(ii), on est ramen\'e au cas o\`u
l'extension r\'esiduelle de~$B/A$ est triviale.
On sait d'ailleurs par \ref{III.1.4}~(i) et~\ref{III.1.3} que les
conditions \'enonc\'ees sont n\'ecessaires. Pour la suffisance,
il suffit de remarquer que la d\'emonstration de~\ref{III.1.5}
prouve, sous les hypoth\`eses faites ici, que $B$ est une
alg\`ebre de s\'eries formelles sur~$A$ (supposant $A$ et $B$
complets, ce qui est loisible).

\begin{remarque}
\label{III.1.8}
Il ne serait pas difficile de d\'evelopper, pour les homomorphismes
formellement lisses, l'analogue de toutes les propri\'et\'es des
morphismes lisses, \'etudi\'ees dans~\ref{II}. Pour les
propri\'et\'es diff\'erentielles, cela demande cependant une
modification de la d\'efinition habituelle des diff\'erentielles
de K\"ahler (\cf I~\ref{I.1}), les produits tensoriels
compl\'et\'es rempla\c cant les produits tensoriels
ordinaires. Nous nous contentons d'\'evoquer ici ces ab\^imes,
ce qui pr\'ec\`ede \'etant suffisant pour notre propos.

Il reste \`a faire le lien entre la lissit\'e formelle, et la
notion de lissit\'e d\'evelopp\'ee dans~\ref{II} (et dont nous
n'avons encore fait nul usage):
\end{remarque}

\begin{proposition}
\label{III.1.9}
Soit $A\to B$ un homomorphisme local, $B$ \'etant localis\'ee
d'une $A$-alg\`ebre de type fini. Pour que $B$ soit lisse sur~$A$,
il faut et il suffit qu'il soit formellement lisse sur~$A$.
\end{proposition}

Utilisant~\ref{III.1.7} et II~\ref{II.2.1}, on est ramen\'e au cas
o\`u $A$ est un corps.

Utilisant~\ref{III.1.4}~(ii) et II~\ref{II.4.13} une extension convenable
$k'$ de~$k$ nous ram\`ene au cas o\`u l'extension r\'esiduelle
pour $B/k$ est triviale. En vertu de~\ref{III.1.5}
(\resp II~\ref{II.5.2}) $B$ est alors lisse sur~$k$ (\resp formellement
\emph{lisse sur} $k$) si et seulement si $B$ est un anneau local
r\'egulier (\resp son compl\'et\'e est une alg\`ebre de
s\'eries formelles sur~$k$). Or il est bien connu que ces deux
conditions sont \'equivalentes (l'extension r\'esiduelle \'etant
triviale).

\section[Propri\'et\'e de rel\`evement]{Propri\'et\'e de rel\`evement caract\'eristique des
ho\-mo\-mor\-phis\-mes formellement lisses}
\label{III.2}

\begin{theoreme}
\label{III.2.1}
Soit $A \to B$ un homomorphisme local d'anneaux locaux d\'efinissant
une extension r\'esiduelle finie. Les conditions suivantes sont
\'equivalentes:
\begin{enumerate}
\item[(i)]
$B$ est formellement lisse sur~$A$
\item[(ii)] Pour tout homomorphisme local $A \to C$, o\`u $C$ est un
anneau local \emph{complet}, tout id\'eal $J$ de~$C$ contenu dans
$\goth{r}(C)$,
et tout $A$-homomorphisme local $B \to C/J$, il existe un
$A$-homomorphisme (n\'ecessairement local) $B \to C$ qui le
rel\`eve.
\item[(iii)] Pour toute $A$-alg\`ebre $C$ (pas n\'ecessairement un
anneau noeth\'erien) tout id\'eal nilpotent $J$ de~$C$, et tout
$A$-homomorphisme $B \to C$ continu (\ie s'annulant sur une puissance
de~$\goth{r}(B)$), il existe un $A$-homomorphisme $B \to C$
(n\'ecessairement continu lui aussi) qui le rel\`eve.
\item[(iv)] M\^eme \'enonc\'e que \textup{(ii)} et \textup{(iii)}, mais $C$
\'etant un anneau local artinien, fini au-dessus de~$A$.
\item[(iv~bis)] Comme \textup{(iv)}, mais $J$ \'etant de plus de carr\'e
nul.
\end{enumerate}
\end{theoreme}

\begin{remarquestar}
Pour la suite de cet expos\'e, nous nous servirons seulement de
l'implication (iv)$\To$(i) ou (iv~bis)$\To$(i);
l'implication directe (i)$\To$(ii) sera prouv\'ee par une
autre m\'ethode au \no suivant lorsque $B$ est localis\'ee d'une
alg\`ebre de type fini sur~$A$. Rappelons que dans la \og bonne\fg
th\'eorie des th\'eor\`emes de Cohen\footnote{\Cf EGA
$0_{\textup{IV}}$ 19.3, 19.8}, la propri\'et\'e (ii) ou (iii)
devient la d\'efinition des homomorphismes formellement lisses,
alors que \ref{III.1.1} devient une propri\'et\'e caract\'eristique
valable seulement dans le cas d'une extension r\'esiduelle finie. On
fera attention que des propri\'et\'es (ii) et (iii) aucune n'est
plus g\'en\'erale que l'autre. On pourrait donner une
propri\'et\'e (\'equivalente) qui les coiffe toutes deux, en
introduisant un anneau lin\'eairement topologis\'e $C$,
\emph{séparé} et \emph{complet}, un id\'eal \emph{fermé
topologiquement nilpotent} de~$C$, et un homomorphisme continu $A \to
C$ (faisant donc de~$C$ une $A$-alg\`ebre topologique); nous
laisserons cette modification au lecteur.
\end{remarquestar}

\subsubsection*{D\'emonstration de~\ref{III.2.1}}
Nous prouverons (i)$\To$(iii)$\To$(ii), puis (iv)$\To$(i). Comme (ii)$\To$(iv) est trivial, et
l'\'equivalence de (iv) et (iv~bis) se voit par une r\'ecurrence
imm\'ediate sur l'entier $n$ tel que $J^n = 0$, cela ach\`evera la
d\'emonstration.

(i)$\To$(iii). Une r\'ecurrence imm\'ediate nous
ram\`ene au cas o\`u $J^2 = 0$. Comme $C$ est fini sur~$A$, il
existe une puissance $\goth{m}^q$ de l'id\'eal maximal de~$A$ qui annule
$C$. Divisant par $\goth{m}^q$, et notant que $B/\goth{m}^qB$ est
encore formellement lisse
sur~$A/\goth{m}^q$ par \ref{III.1.4}~(i), on peut supposer $A$
artinien. Comme $B$ est plat sur~$A$ par~\ref{III.1.3}, $B$ \emph{est
libre sur} $A$ puisque $A$ est artinien. Donc il existe un
\emph{homomorphisme de} $A$-\emph{modules}
$$
w \colon B \to C
$$
qui rel\`eve l'homomorphisme donn\'e $u \colon B \to C/J$. Posons
$$
f(x,y) = w(xy) - w(x)w(y) \qquad (x,y \in B)
$$
alors $f(x,y) \in J$, et $f$ est donc une application
$A$-bilin\'eaire de~$B \times B$ dans $J$. Pour qu'il existe un
rel\`evement $v \colon B \to C$ de~$u$ qui soit un homomorphisme
d'alg\`ebres, il faut et il suffit qu'il existe une application
$A$-lin\'eaire $g \colon B \to J$ telle que $v = w + g$ soit un
homomorphisme d'alg\`ebres, ce qui s'\'ecrit
\begin{align*}
g(1) &= 1 - w(1) \\
g(x,y) - u(x)g(y) -u(y)g(x) &= -f(x,y) \qquad (\text{pour } x,y \in
B)
\end{align*}
C'est l\`a un syst\`eme d'\'equations \emph{linéaires} dans
$\Hom_A(B,J)$, \`a seconds membres dans $J$, donc il a une solution
si et seulement si le syst\`eme correspondant dans $\Hom_A(B,J)
\otimes_A A'$, \`a seconds membres dans $J' = J \otimes_A A'$, a une
solution, --- $A'$ d\'esignant une alg\`ebre fid\`element plate
sur~$A$. Or soit $A'$ une alg\`ebre finie et libre sur~$A$, locale,
telle que $B' = B \otimes_A A'$ soit une alg\`ebre de s\'eries
formelles sur~$A'$ (N.B. on peut dans notre d\'emonstration supposer
$A$ et $B$ complets, comme on constate aussit\^ot). Comme $A'$ est
libre de type fini sur~$A$, on a
$$
\Hom_A(B,J) \otimes_A A' = \Hom_{A'}(B',J')
$$
et on constate que le syst\`eme d'\'equations obtenu dans
$\Hom_{A'}(B',J')$ est celui qui d\'etermine les homomorphismes de
$A'$-alg\`ebres $B' \to C' = C \otimes_A A'$ qui rel\`event
l'homomorphisme $u' \colon B' \to C'/J'$ d\'eduit de~$u$ par
extension des scalaires en \og corrigeant\fg par un homomorphisme de
$A'$-modules $g' \colon B' \to J'$ l'homomorphisme de~$A'$-modules $w'
\colon B' \to C'$ d\'eduit de~$w$ par extension des
scalaires. (Noter que $B$ engendre $B'$ comme $A'$-module). On est
ainsi ramen\'e \`a prouver (iii) lorsque $B$ est une
\emph{algèbre de séries formelles} sur~$A$, $B = A[[t_1,
\dots, t_n]]$. Relevons alors de fa\c con quelconque les images dans
$C/J$ des $t_i$ en des \'el\'ements $z_i$ de~$C$. Comme les $z_i$~mod~$J$ sont nilpotents ($u \colon B \to C/J$ \'etant continu) il en
est de m\^eme des $z_i$ (puisque $J$ est nilpotent), donc les $z_i$
d\'efinissent un homomorphisme continu de~$A$-alg\`ebres
topologiques de~$B$ dans $C$ discret, relevant \'evidemment $u$,
cqfd.


(iii)$\To$(ii).
Soit $\goth{n}$ l'id\'eal maximal de~$C$, et
pour tout entier $q > 0$, soit
$$
C_q = C/\goth{n}^q \text{\quoi, \quad} J_q = (J +
\goth{n}^q)/\goth{n}^q
$$
donc $C_q/J_q$ s'identifie \`a une alg\`ebre quotient de~$C/J$,
d'autre part l'homomorphisme compos\'e $u_q \colon B \to C/J \to
C_q/J_q$ est continu de~$B$ dans $C_q/J_q$ discret, et $J_q$ est un
id\'eal nilpotent dans $C_q$. On construit alors de proche en proche
des $A$-homomorphismes
$$
v_q \colon B \to C_q
$$
tels que (a) $v_q$ rel\`eve $u_q$ et (b) $v_q$ rel\`eve
$v_{q-1}$. La possibilit\'e de la r\'ecurrence se v\'erifie
ais\'ement, car comme
$$
u_q \colon B \to C/(J + \goth{n}^q) \text{\quad et \quad} v_{q-1}
\colon B \to C/\goth{n}^{q-1}
$$
d\'efinissent le m\^eme homomorphisme
$$
B \to C/((J + \goth{n}^q) + \goth{n}^{q-1}) = C/(J + \goth{n}^{q-1}) =
C_{q-1}/J_{q-1}
$$
\`a savoir $u_{q-1}$, ils d\'efinissent un homomorphisme
$$
B \to C/J_q' \text{\quad où \quad} J_q' = (J + \goth{n}^q)\cap
\goth{n}^{q-1} \supset \goth{n}^q
$$
(dont ils proviennent l'un et l'autre par r\'eduction). On est donc
ramen\'e \`a relever un homomorphisme $B \to C/J_q'$ de~$B$ dans
un quotient de~$C_q$ par un id\'eal $J_q'/\goth{n}^q$ contenu dans
$J_q$, donc nilpotent, et cela est possible d'apr\`es
l'hypoth\`ese (iii).

Ceci fait, les $v_q$ d\'efinissent un homomorphisme de~$B$ dans la
limite projective $C$ des $C_q$. Comme $J$ est ferm\'e, $J$ est la
limite projective des $J_q$, donc $v$ rel\`eve $u$, cqfd.

(iv)$\To$(i). On constate d'abord aussit\^ot que si (iv)
est v\'erifi\'e, (iv) reste v\'erifi\'e pour las composants
locaux de~$B \otimes_A A'$ sur~$A'$, si $A'$ est une alg\`ebre
locale finie sur~$A$. Prenant $A'$ libre sur~$A$ et telle que les
extensions r\'esiduelles de~$B'$ au-dessus de~$A'$ soient triviales,
on est ramen\'e gr\^ace \`a~\ref{III.1.4} (ii) au cas o\`u
l'extension r\'esiduelle de~$B$ sur~$A$ est triviale. Nous allons
prouver alors le r\'esultat un peu pr\'ecis:

\begin{corollaire}
\label{III.2.2}
Sous les conditions de~\ref{III.2.1}, supposons de plus l'extension
r\'esiduelle de~$B$ au-dessus de~$A$ triviale. Alors les conditions
\'equivalentes
de~\ref{III.2.1} \'equivalent aussi aux deux conditions suivantes
(en supposant dans \textup{(v)} $A$ et $B$ complets):
\begin{enumerate}
\item[(iv~ter)] Comme \textup{(iv)}, mais l'anneau local artinien $C$ fini sur~$A$ \'etant restreint \`a avoir une extension r\'esiduelle
triviale (et de plus, si on y tient, l'id\'eal $J$ \'etant de
carr\'e nul).
\item[(v)] Il existe un $A$-homomorphisme local (o\`u $n =
\dim{\goth{n}/(\goth{n}^2+\goth{m}B})$)
$$
u \colon B \to B_1 = A[[t_1, \dots, t_n]]
$$
induisant un \emph{isomorphisme}
$$
\goth{n}/(\goth{n}^2+\goth{m}B) \isomto
\goth{n}_1/(\goth{n}_1^2+\goth{m}B_1)
$$
o\`u $\goth{n}$ ($\goth{n}_1$) est l'id\'eal maximal de~$B$
($B_1$), $\goth{m}$ celui de~$A$.
\end{enumerate}
\end{corollaire}

\subsubsection*{D\'emonstration} Comme (iv~bis) implique \'evidemment (iv~ter) --- en faisant abstraction du canular de l'id\'eal de carr\'e nul ---,
il suffira de prouver (iv~ter)$\To$(v)$\To$(i).

(iv~ter)$\To$(v). Choisissons une base $a_1,\dots,a_n$ de
$\goth{n}/(\goth{n}^2+\goth{m}B)$, ce qui d\'efinit donc un
homomorphisme local de~$A$-alg\`ebres
$$
B \to B_1/(\goth{n}_1^2+\goth{m}B_1) = k[t_1, \dots, t_n]/(t_1, \dots,
t_n)^2
$$
qu'on peut relever de proche en proche, en vertu de (iv~ter) en des
homomorphismes de~$A$-alg\`ebres de~$B$ dans $B_1/\goth{n}_1^2$,
$B_1/\goth{n}_1^3$ etc, d'o\`u en passant \`a la limite projective
l'homomorphisme $B \to B_1$ ayant la propri\'et\'e voulue.

(v)$\To$(i). Comme dans le diagramme commutatif
$$
\xymatrix{{\goth{m}/\goth{m}^2}\ar[r]\ar[d]&{\goth{n}/\goth{n}^2}\ar[r]\ar[d]&{\goth{n}/(\goth{n}^2+\goth{m}B)}\ar[r]\ar[d]&{0}\\{\goth{m}/\goth{m}^2}\ar[r]&{\goth{n}_1/\goth{n}_1^2}\ar[r]&{\goth{n}_1/(\goth{n}_1^2+\goth{m}B_1)}\ar[r]&{0}}
$$
les deux lignes sont exactes, et les fl\`eches verticales
extr\^emes surjectives, la fl\`eche m\'ediane est surjective et
il s'ensuit ($B$ \'etant complet) que $B \to B_1$ est
\emph{surjectif}. Soient $x_i$ ($1 \leq i \leq n$) des
\'el\'ements de~$B$ qui rel\`event les $t_i$. Ils
d\'efinissent donc un homomorphisme de~$A$-alg\`ebre $B_1 \to B$,
qui sera surjectif pour la m\^eme raison que $u$, et dont le
compos\'e avec $u$ est l'identit\'e par construction. Donc $B_1
\to B$ est aussi injectif, et est par suite un isomorphisme. On trouve
donc

\begin{corollaire}
\label{III.2.3}
Sous
les conditions de \ref{III.2.2}~\textup{(v)}, $u$ est n\'ecessairement
un isomorphisme.
\end{corollaire}

Cela ach\`eve de prouver que $B$ est formellement lisse sur~$A$. On
a d'ailleurs en m\^eme temps retrouv\'e~\ref{III.1.5} (mais il n'y
a gu\`ere de m\'erite \`a \c ca).

\section[Prolongement infinit\'esimal local des morphismes]{Prolongement infinit\'esimal local des morphismes dans un
$S$-sch\'ema lisse}
\label{III.3}

\begin{theoreme}
\label{III.3.1}
Soit $f \colon X \to Y$ un morphisme localement de type fini.
Conditions \'equivalentes:
\begin{enumerate}
\item[(i)] $f$ est lisse.
\item[(ii)] Pour tout pr\'esch\'ema $Y'$ sur~$Y$, tout
sous-pr\'esch\'ema ferm\'e $Y'_0$ de~$Y'$ ayant m\^eme espace
sous-jacent que $Y'$, tout $Y$-morphisme $g_0 \colon Y'_0 \to X$ et
tout $z \in Y'_0$, il existe un voisinage ouvert $U$ de~$z$ dans $Y'$
et un prolongement $g$ de~$g_0|Y'_0 \cap U$ en un $Y$-morphisme $U \to
X$.
\item[(ii~bis)] Pour $Y'$, $Y'_0$ et $z$ comme dans \textup{(ii)}, posant
$X' = X \times_Y Y'$, $X'_0 = X \times_Y Y'_0$, toute section de
$X'_0$ sur~$Y'_0$ se prolonge en une section de~$X'$ au-dessus d'un
voisinage ouvert $U$ de~$z$.
\item[(iii)] Pour tout $Y$-sch\'ema $Y'$, spectre d'un anneau artinien
local fini sur quelque $\mathcal{O}_y$ ($y \in Y$), tout
sous-pr\'esch\'ema ferm\'e non vide~$Y'_0$ de~$Y'$, et tout
$Y$-morphisme $g_0 \colon Y'_0 \to X$, il existe un $Y$-morphisme $g
\colon Y' \to X$ qui prolonge $g_0$.
\item[(iii~bis)] Pour tout $Y'$, $Y'_0$ comme dans \textup{(iii)},
posant $X' = X \times_Y Y'$, $X'_0 = X \times_Y Y'_0$, toute section
de~$X'_0$ sur~$Y'_0$ se prolonge en une section de~$X'$ sur~$Y'$.
\end{enumerate}
\end{theoreme}

\subsubsection*{D\'emonstration} L'\'equivalence de (ii) et (ii~bis) d'une part, de
(iii) et (iii~bis) d'autre part, est triviale, ainsi que l'implication
(ii)$\To$(iii). Il reste donc \`a prouver (i)$\To$(ii) et (iii)$\To$(i).

(i)$\To$(ii). Soit $x = g_0(z)$. Rempla\c cant $X$ par un
voisinage ouvert convenable de~$x$, et $Y'$ par le pr\'esch\'ema
induit sur l'ouvert image r\'eciproque de ce dernier par~$g_0$, on
peut supposer que $X$ est \'etale au-dessus de~$Y[t_1, \dots, t_n]$.
Consid\'erons le $Y$-morphisme compos\'e $Y'_0 \to X \to Y[t_1,
\dots, t_n]$, il est d\'efini par $n$ sections du faisceau
$\mathcal{O}_{Y'_0}$, qui peuvent donc se prolonger au voisinage de
$z$ en des sections de~$\mathcal{O}_{Y'}$, donc on peut supposer que
le morphisme en question a \'et\'e prolong\'e en un
$Y$-morphisme $Y' \to Y'_0$. En vertu de (I~\ref{I.5.6}) il existe
alors un unique $Y$-morphisme $g \colon Y' \to X$ qui rel\`eve le
pr\'ec\'edent, et prolonge en m\^eme temps $g_0$, cqfd.

(iii)$\To$(i). Comme l'ensemble des points o\`u $f$ est
lisse est ouvert, il suffit de prouver qu'il contient tout $x \in X$
qui est \emph{fermé} dans sa fibre. Soit $y = f(x)$, alors
$\mathcal{O}_x$ est une alg\`ebre sur~$\mathcal{O}_y$, localis\'ee
d'une alg\`ebre de type fini, \`a extension r\'esiduelle finie.
D'autre part, l'hypoth\`ese (iii) implique que tout homomorphisme de
$\mathcal{O}_x$ dans une alg\`ebre $A/J$, o\`u $A$ est une
alg\`ebre artinienne locale finie sur~$\mathcal{O}_y$, et $J$ un
id\'eal contenu dans son radical, se rel\`eve en un homomorphisme
de~$\mathcal{O}_x$ dans l'alg\`ebre $A$ (compte tenu qu'un morphisme
d'un $\Spec(B)$, $B$ anneau local, dans $X$ est d\'etermin\'e
biunivoquement par un homomorphisme local d'un $\mathcal{O}_x$ ($x \in
X$) dans $B$). Par~\ref{III.2.1} il s'ensuit que $\mathcal{O}_x$ est
formellement lisse sur~$\mathcal{O}_y$, donc lisse sur~$\mathcal{O}_y$
en vertu de~\ref{III.1.9}.

\begin{corollaire}
\label{III.3.2}
Soit $f\colon X\to Y$ comme dans~\ref{III.3.1}. Conditions
\'equivalentes:
\begin{enumerate}
\item[(i)] $f$ est \'etale
\item[(ii)] On a la condition~\textup{(ii)} de~\ref{III.3.1} avec
{\emph{unicité}} du prolongement $g$ de~$g_0$ \`a $U$.
\item[(iii)] On a la condition~\textup{(iii)} de~\ref{III.3.1} avec
{\emph{unicité}} de~$g$.
\end{enumerate}
\end{corollaire}

Il suffit de noter, dans la d\'emonstration de (i)$\To$(ii)
ci-dessus, qu'on ne peut avoir unicit\'e (quand $Y'_0$ n'est pas
identique \`a~$Y'$ au voisinage de~$z$) que si $n=0$ (condition
qu'on sait \^etre suffisante).

\begin{corollaire}
\label{III.3.3}
Soit $X$ un pr\'esch\'ema localement de type fini sur un anneau
local {\emph{complet}} $A$, $y$ le point ferm\'e de~$Y=\Spec(A)$ et
$x$ un point de~$f^{-1}(y)$ {\emph{rationnel}}
sur~$\kres(y)$. Si $X$
est {\emph{lisse sur~$A$ en~$x$}}, alors il existe une section $s$ de
$X$ sur~$Y$ \og passant par~$x$\fg \ie telle que $s(y)=x$.
\end{corollaire}

En particulier, si $X$ est lisse sur~$A$, alors l'application
naturelle
$$
\Gamma(X/Y)\to\Gamma(X\otimes_A k/k)
$$
des sections de~$X$ sur~$Y$ dans l'ensemble des points de la fibre
$f^{-1}(y)=X\otimes_A k$ rationnels sur~$k$, est surjective. Ce fait
\'etait surtout bien connu et utilis\'e lorsque $A$ est un anneau
de valuation discr\`ete et $X$
est propre sur~$A$ (en fait, projectif sur~$A$), auquel cas les
sections de~$X$ sur~$Y$ (\ie les \og points de~$X$ \`a valeurs
dans~$A$\fg) s'identifient aussi aux sections rationnelles, \ie aux
points de~$X\otimes_A K=X_K$ (qui est un sch\'ema propre et simple
sur~$K$) \`a valeurs dans $K=$ corps des fractions
de~$A$,
\ie aux
points de~$X$ rationnels sur~$K$.

\section{Prolongement infinit\'esimal local des $S$-sch\'emas lisses}
\label{III.4}

\begin{theoreme}
\label{III.4.1}
Soient $Y$ un pr\'esch\'ema localement noeth\'erien, $Y_0$ un
sous-pr\'e\-sch\'ema ferm\'e ayant m\^eme espace sous-jacent,
$X_0$ un $Y_0$-pr\'esch\'ema lisse, $x$ un point de~$X_0$. Alors
il existe un voisinage ouvert $U_0$ de~$x$, un pr\'esch\'ema $X$
lisse sur~$Y$, et un $Y_0$-isomorphisme:
$$
h\colon U_0\isomto X\times_Y Y_0.
$$
De plus, si $(U'_0,X',h')$ est une autre solution de ce probl\`eme,
alors \og elle est isomorphe \`a la premi\`ere au voisinage
de~$x$\fg.
\end{theoreme}

On laisse au lecteur de pr\'eciser ce qu'on veut dire par l\`a. On
peut noter que pour $U_0$ donn\'e, une solution du
probl\`eme pos\'e
revient \`a la donn\'ee, sur~$U_0$, d'un faisceau d'alg\`ebres
$\cal{B}$ sur~$f_0^{-1}(\cal{O}_Y)$ (o\`u $f_0$ est l'application
continue sous-jacente au morphisme structural $U_0\to Y_0$), muni d'un
homomorphisme d'anneaux $\cal{B}\to \cal{O}_{U_0}$ compatible avec
l'homomorphisme $f^{-1}(\cal{O}_Y)\to f^{-1}(\cal{O}_{Y_0})$, tels que

(a) cet homomorphisme induit un isomorphisme
$$
\cal{B}\otimes_{f^{-1}(\cal{O}_Y)} f^{-1}(\cal{O}_{Y_0})\isomto
\cal{O}_{Y_0};
$$

(b) $U_0$ muni de~$\cal{B}$ devient un $Y$-pr\'esch\'ema lisse.

De cette fa\c con, le sens pr\'ecis de l'assertion d'unicit\'e
locale devient particuli\`erement \'evident.

\subsubsection*{D\'emonstration} On peut supposer d\'ej\`a que $X_0$ est
\'etale \hbox{sur un $Y_0[t_1,\dots,t_n]=Y'_0$}. Or ce dernier peut
\^etre consid\'er\'e comme un sous-pr\'esch\'ema ferm\'e
de~$Y'{=} Y[t_1,\dots,t_n]$, ayant m\^eme espace sous-jacent. Par
(I~\ref{I.8.3}), il existe un $X$ \'etale sur~$Y'$, et un
$Y'_0$-isomorphisme $X\times_{Y'}Y'_0\isomto X'$. On a gagn\'e pour
l'existence. Pour l'unicit\'e, on utilise la
propri\'et\'e~\ref{III.3.1}~(ii)
des morphismes lisses, en tenant compte du

\begin{lemme}
\label{III.4.2}
Soient $Y$ un pr\'esch\'ema, $Y_0$ un sous-pr\'esch\'ema
ferm\'e d\'efini par un faisceau d'id\'eaux $\cal{J}$ localement
nilpotent, $X$ et $X'$ des $Y$-pr\'esch\'emas, $u\colon X\to X'$
un $Y$-morphisme. On suppose $X$ plat sur~$Y$. Pour que $u$ soit un
isomorphisme, il faut et suffit que $u_0\colon X\times_Y Y_0\to
X'\times_Y Y_0$ le soit.
\end{lemme}

D\'emonstration facile, en passant au cas affine et regardant les
gradu\'es associ\'es. On notera d'ailleurs que l'\'enonc\'e
analogue, obtenu en rempla\c cant \og isomorphisme\fg par \og immersion
ferm\'ee\fg, est \'egalement valable, et sans hypoth\`ese de
platitude.

\begin{remarque}
\label{III.4.3}
Il est essentiel de noter que le prolongement local $X$ obtenu
dans~\ref{III.4.1} {\emph{n'est pas canonique}}, en d'autres termes
l'isomorphisme local entre deux solutions n'est pas unique, \ie il
existe en g\'en\'eral des $Y$-automorphismes non triviaux de~$X$
induisant l'identit\'e sur le sous-pr\'esch\'ema ferm\'e
$X_0=X\times_Y Y_0$. C'est pour cela qu'il faut s'attendre, pour la
construction de prolongements infinit\'esimaux {\emph{globaux}} de
pr\'esch\'emas
lisses, \`a l'existence d'une obstruction de
nature cohomologique, qui sera pr\'ecis\'ee plus bas (\no \ref{III.6}).
\end{remarque}

\section{Prolongement infinit\'esimal global des morphismes}
\label{III.5}

Soient $T$ un espace topologique, $\cal{G}$ un faisceau de groupes
sur~$X$,
$\cal{P}$
un faisceau d'ensembles sur~$T$ sur lequel $\cal{G}$ op\`ere (\`a
droite, pour fixer les id\'ees). On dit que $\cal{P}$ est
\emph{formellement principal homogène}
\index{formellement principal homog\`ene (faisceau)|hyperpage}sous~$\cal{G}$, si l'homomorphisme bien connu
$$
\cal{G}\times \cal{P}\to \cal{P}\times \cal{P}
$$
de faisceaux d'ensembles, d\'eduit des op\'erations de~$\cal{G}$
sur~$\cal{P}$, est un {\emph{isomorphisme}}. Il revient au m\^eme de
dire que pour tout $x\in T$, $\cal{P}_x$ est {\emph{vide ou un espace
principal homogène}} sous le groupe ordinaire $\cal{G}_x$, ou
aussi que pour tout ouvert $U$ de~$T$, $\cal{P}(U)$ est vide ou un
espace principal homog\`ene sous le groupe ordinaire
$\cal{G}(U)$. On dit que
$\cal{P}$
est un {\emph{faisceau principal homogène}}
\index{principal homog\`ene (faisceau)|hyperpage}sous~$\cal{G}$
s'il l'est formellement et si en plus les $\cal{P}_x$ sont non vides
(en d'autres termes, si {\emph{tous}} les $\cal{P}_x$ sont des espaces
principaux homog\`enes,
donc non vides, sous les $\cal{G}_x$)\footnote{Il semble
pr\'ef\'erable d'adopter le terme plus court et plus parlant de
\og {\emph{torseur sous}}~$\cal{G}$\fg,
introduit dans la th\`ese de J.\, \textsc{Giraud}.}. Rappelons que l'ensemble
des classes (\`a isomorphisme pr\`es) de faisceaux principaux
homog\`enes sous~$\cal{G}$ s'identifie \`a l'ensemble de
cohomologie $\H^1(T,\cal{G})$, qui est aussi le groupe de cohomologie
usuel de~$T$ \`a coefficients dans $\cal{G}$ lorsque $\cal{G}$ est
commutatif. On a ainsi, pour tout $\cal{P}$ principal homog\`ene,
une classe caract\'eristique $c(\cal{P})\in\H^1(T,\cal{G})$, dont la
trivialit\'e est n\'ecessaire et suffisante pour que $\cal{P}$
soit trivial (\ie isomorphe \`a~$\cal{G}$, sur lequel $\cal{G}$
op\`ere par translations \`a droite), ou encore pour que $\cal{P}$
ait une section.

\begin{proposition}
\label{III.5.1}
Soient $S$ un pr\'esch\'ema, $X$ et $Y$ des pr\'esch\'emas
sur~$S$, $Y_0$ un sous-pr\'esch\'ema ferm\'e de~$Y$ d\'efini
par un Id\'eal $\cal{J}$ sur~$Y$ {\emph{de carré nul}}. Soit
$g_0$ un $S$-morphisme de~$Y_0$ dans~$X$, et $\cal{P}(g_0)$ le
faisceau sur~$Y$ dont les sections sur un ouvert~$U$ sont les
prolongements $g\colon U\to X$ de~$g_0|U\cap Y_0$ en un $S$-morphisme
$g$. Alors $\cal{P}(g_0)$ est un faisceau {\emph{formellement
principal homogène}} (de fa\c con naturelle) sous le faisceau
en groupes commutatif
$$
\cal{G}=\SheafHom_{\cal{O}_{Y_0}}(g_0^*(\it{\Omega}^1_{X/S}),\cal{J})
$$
\end{proposition}

Posons $\cal{P}=\cal{P}(g_0)$. Nous devons d\'efinir pour tout
ouvert $U$ de~$Y$ une application
$$
\cal{P}(U)\times \cal{G}(U)\to \cal{P}(U)
$$
de fa\c con que (a) pour $g\in \cal{P}(U)$ fix\'e, l'application
$s\mto gs$ de~$\cal{G}(U)$ dans~$\cal{P}(U)$ est bijective (b)
$\cal{P}(U)$ devient un ensemble \`a groupe d'op\'erateurs
$\cal{G}(U)$ (c) les applications pr\'ec\'edentes sont compatibles
avec les op\'erateurs de restriction pour un ouvert $V\subset U$. La
v\'erification de (c) sera triviale, on peut donc pour simplifier
supposer $U=Y$. La v\'erification de (b) (qui est, si on veut, de
nature locale) sera laiss\'ee au lecteur, nous nous bornerons donc,
pour un $g\in \cal{P}(Y)$ donn\'e, de d\'efinir une bijection
naturelle de~$\cal{G}(Y)$ sur~$\cal{P}(Y)$. Donc on suppose donn\'e
d\'ej\`a un $S$-morphisme $g\colon Y\to X$, et on cherche une
bijection canonique
\begin{equation*}
\label{eq:III.5.1.*}
\tag{$*$}
\Hom_{\cal{O}_{Y_0}}(g_0^*(\it{\Omega}^1_{X/S}),\cal{J})\isomto
\cal{P}(Y)
\end{equation*}
o\`u $\cal{P}(Y)$ est l'ensemble des $S$-morphismes $g'$ de~$Y$
dans~$X$ induisant le m\^eme morphisme $g_0\colon Y_0\to X$ que~$g$.
La donn\'ee d'un tel $g'$ est \'equivalente \`a la donn\'ee
d'un $S$-morphisme $h\colon Y\to X\times_S X$ tel que
${\pr}_1\circ h=g$, et $h\circ i=(g_0,g_0)$, o\`u
${\pr}_1\colon X\times_S X \to X$
est la premi\`ere projection, $i\colon Y_0\to Y$ l'immersion
canonique, et $(g_0,g_0)\colon Y_0\to X\times_S X$ est le morphisme
$\Delta_{X/S}g_0$ de
composantes
$g_0,g_0$:
$$
\xymatrix@C=2.2cm{{X\times_S X}\ar[d]_{{\pr}_1}&{Y_0}\ar[l]_-{h_0=(g_0,g_0)}\ar[d]^i\\{X}&{Y}\ar[l]_-g\ar@{-->}[lu]_-h}
$$
Comme $h_0$ se factorise par l'immersion diagonale $\Delta_{X/S}$ et
que $Y$ est dans le voisinage infinit\'esimal d'ordre $1$ de~$Y_0$
(\ie $\cal{J}^2=0$) les $h$ cherch\'es se factorisent
n\'ecessairement (de fa\c con unique) par le voisinage
infinit\'esimal du premier ordre de la diagonale, lequel s'identifie
(en tant que $X$-pr\'esch\'ema gr\^ace \`a ${\pr}_1$)
au spectre $X'$ du faisceau d'alg\`ebres
$\cal{O}_X+\it{\Omega}^1_{X/S}$ (o\`u le deuxi\`eme terme est
consid\'er\'e comme un Id\'eal de carr\'e nul), le morphisme
diagonal $X\to X'$ correspondant \`a l'augmentation canonique de ce
faisceau d'alg\`ebres. Posons $Y'=X'\times_X Y$, $Y'_0=Y'\times_Y
Y_0=X'\times_X Y_0$, de sorte que les $h$ cherch\'es sont en
correspondance biunivoque avec les sections $u$ de~$Y'$ sur~$Y$ qui
prolongent une section donn\'ee $u_0$ de~$Y'_0$ sur~$Y_0$. On peut
d'ailleurs identifier $Y'$ au spectre du faisceau d'alg\`ebres sur~$Y$ $\cal{A}=g^*(\cal{O}_X+\it{\Omega}^1_{X/S})=\cal{O}_Y+
g^*(\it{\Omega}^1_{X/S})$, et $Y'_0$ au faisceau d'alg\`ebres
$\cal{A}_0=\cal{A}\otimes_{\cal{O}_Y}\cal{O}_{Y_0}=\cal{O}_{Y_0}+
g^*_0(\it{\Omega}^1_{X/S})$, alors la section $u_0$ est celle
d\'efinie par l'augmentation canonique de~$A_0$
dans~$\cal{O}_{Y_0}$. Donc $\cal{P}(Y)$ s'identifie \`a l'ensemble
des homomorphismes d'alg\`ebres $\cal{A}\to \cal{O}_Y$ qui induisent
l'augmentation canonique $\cal{A}_0\to \cal{O}_{Y_0}$. Or les
homomorphismes d'alg\`ebres $\cal{A}\to \cal{O}_Y$
correspondent
biunivoquement aux homomorphismes de Modules $\cal{M}\to \cal{A}$
(posant pour simplifier $\cal{M}=g^*(\it{\Omega}^1_{X/S})$), et on
s'int\'eresse \`a ceux qui induisent l'homomorphisme {\emph{nul}}
$\cal{M}_0\to \cal{O}_{Y_0}$ (o\`u
$\cal{M}_0=\cal{M}\otimes_{\cal{O}_Y}\cal{O}_{Y_0}$) \ie qui
appliquent $\cal{M}$ dans l'id\'eal $\cal{J}$ de l'augmentation. On
trouve donc l'ensemble
$\Hom_{\cal{O}_Y}(\cal{M},\cal{J})=\Hom_{\cal{O}_{Y_0}}(\cal{M}_0,\cal{J})$
(puisque $\cal{J}$ est annul\'e par $\cal{J}$). C'est la bijection
canonique cherch\'ee \eqref{eq:III.5.1.*}.

Tenant compte de l'implication (i)$\To$(iii)
dans~\ref{III.3.1}, on trouve:

\begin{corollaire}
\label{III.5.2}
Si $X$ est lisse sur~$S$ (du moins aux points de~$g_0(Y_0)$)
alors~$\cal{P}$ est m\^eme un \emph{faisceau principal homogène}
sous le faisceau en groupes commutatifs~$\cal{G}$, qui en l'occurrence
peut aussi s'\'ecrire
\begin{equation*}
\cal{G}=g_0^*(\goth{g}_{X/S})\otimes_{\cal{O}_{Y_0}}\cal{J}
\end{equation*}
o\`u $\goth{g}_{X/S}$
\label{indnot:cb}est le faisceau sur~$X$ dual de~$\it{\Omega}_{X/S}^1$, \ie le
\emph{faisceau tangent}
\index{faisceau tangent|hyperpage}\index{tangent (faisceau)|hyperpage}(ou \emph{faisceau des dérivations})
\index{derivations (faisceau des)@d\'erivations (faisceau des)|hyperpage}\index{faisceau des derivations@faisceau des d\'erivations|hyperpage}de~$X$ par rapport \`a~$S$. (Cette derni\`ere formule provient du
fait que~$\it{\Omega}_{X/S}^1$ est alors libre de type fini).
\end{corollaire}

En particulier, \`a ce faisceau principal homog\`ene correspond
une classe de cohomologie dans $\H^1(Y_0,\cal{G})$, dont l'annulation
est n\'ecessaire et suffisante pour l'existence d'un
$S$-morphisme~$g$ prolongeant~$g_0$. Et s'il existe un tel
prolongement, l'ensemble de tous les prolongements possibles est un
espace homog\`ene sous le groupe~$\H^0(Y_0,\cal{G})$.

Dans l'application des m\'ethodes de la g\'eom\'etrie formelle,
la situation est le plus souvent la suivante: on donne deux
$S$-pr\'esch\'emas $X$ et~$Y$, un Id\'eal coh\'erent~$\cal{I}$
sur~$S$, on d\'esigne par~$S_n$ le sous-pr\'esch\'ema ferm\'e
de~$S$ d\'efini par~$\cal{I}^{n+1}$, et on pose
\begin{equation*}
X_n=X\times_S S_n \quoi,\quad Y_n=Y\times_S S_n.
\end{equation*}
On suppose qu'on a un $S_n$-morphisme:
\begin{equation*}
g_n\colon Y_n\to X_n
\end{equation*}
(ou, ce qui revient au m\^eme, un~$S$-morphisme $Y_n\to X$ ou encore
un $S_{n+1}$-morphisme $Y_n\to X_{n+1}$,
puisqu'un
tel morphisme
induit n\'ecessairement $Y_n\to X_n$), et on cherche \`a le
prolonger en un $S_{n+1}$-morphisme
\begin{equation*}
g_{n+1}\colon Y_{n+1}\to X_{n+1}
\end{equation*}
(Si on peut continuer ind\'efiniment, on obtient donc un morphisme
$\widehat{Y}\to\widehat{X}$ pour les pr\'esch\'emas formels
compl\'et\'es de~$Y$ et~$X$ pour les Id\'eaux
$\cal{IO}_Y$ et~$\cal{IO}_X$).
On peut appliquer~\ref{III.5.1} en y rempla\c cant $(S,X,Y,Y_0,g_0)$
par $(S_{n+1},X_{n+1},Y_{n+1},Y_n,g_n)$, le faisceau $\cal{G}$ devient
ici le faisceau des homomorphismes de Modules
de~$g_n^*(\it{\Omega}_{X_{n+1}/S_{n+1}}^1)$
dans~$\cal{J}=\cal{I}^{n+1}\cal{O}_Y/\cal{I}^{n+2}\cal{O}_Y$.
Comme~$\cal{J}$ est annul\'e par $\cal{IO}_Y$, on peut remplacer
alors~$g_n^*(\it{\Omega}_{X_{n+1}/S_{n+1}}^1)$
par le faisceau qu'il induit sur~$Y_0$,
savoir~$h_0^*(\it{\Omega}_{X/S}^1)$, o\`u~$h_0$ est le compos\'e
$Y_0\to Y_n\to X_{n+1}$, ou encore le compos\'e $Y_0 \to X_0\to
X_{n+1}$, o\`u $g_0\colon Y_0\to X_0$ est induit par $g_n$. Comme
l'image inverse de~$\it{\Omega}_{X_{n+1}/S_{n+1}}^1$
sur~$X_0=X_{n+1}\times_{S_{n+1}}S_0$
est~$\it{\Omega}_{X_{0}/S_{0}}^1$, on voit qu'on a aussi
\begin{equation*}
\cal{G}=\SheafHom_{\cal{O}_{Y_0}} (g_0^*(\it{\Omega}_{X_0/S_0}^1),
\cal{I}^{n+1}\cal{O}_Y/\cal{I}^{n+2}\cal{O}_Y)
\end{equation*}
Donc on obtient le

\begin{corollaire}
\label{III.5.3}
Soient $S,X,Y,\cal{I},g_n$ comme ci-dessus, soit $\cal{P}(g_n)$ le
faisceau sur~$Y$ dont les sections sur un ouvert $U$ sont les
prolongements $g_{n+1}$ de~$g_n$ en un $S_{n+1}$-morphisme $Y_{n+1}\to
X_{n+1}$. Alors
$\cal{P}(g_{n})$
est un faisceau formellement principal homog\`ene sous le faisceau
en groupes
\begin{equation*}
\cal{G}=\SheafHom_{\cal{O}_{Y_0}} (g_0^*(\it{\Omega}_{X_0/S_0}^1),
\gr_{\cal{IO}_Y}^{n+1}(\cal{O}_Y))
\end{equation*}
\end{corollaire}

En particulier:

\begin{corollaire}
\label{III.5.4}
Si de plus $X$ est lisse sur~$S$ (du moins aux points de~$g_0(Y_0)$)
alors
$\cal{P}(g_{n})$
est m\^eme un faisceau principal homog\`ene.
\enlargethispage{.5cm}En particulier, il d\'efinit une classe d'obstruction dans
$\H^1(Y_0,\cal{G})$, dont l'annulation est n\'ecessaire et
suffisante pour l'existence d'un prolongement global~$g_{n+1}$
de~$g_n$. Et s'il existe un tel prolongement, l'ensemble de tous les
prolongements globaux est un espace principal homog\`ene sous
$\H^0(Y_0,\cal{G})$. Enfin, dans le cas envisag\'e, le
faisceau~$\cal{G}$ peut aussi s'\'ecrire
\begin{equation*}
\cal{G}=g_0^*(\goth{g}_{X_0/S_0})\otimes_{\cal{O}_{Y_0}}
\gr_{\cal{IO}_Y}^{n+1}(\cal{O}_Y)
\end{equation*}
\end{corollaire}

Proc\'edant de proche en proche, on voit donc que si tous les
$\H^1(Y_0,\cal{G}_n)$ son nuls (o\`u
$\cal{G}_n=g_0^*(\goth{g}_{X_0/S_0})\otimes
\gr_{\cal{IO}_Y}^{n}(\cal{O}_Y)$), alors partant avec un $g_k$
quelconque, on peut le prolonger successivement en $g_{k+1},\dots$ En
particulier, si $\cal{I}$ est nilpotent, on pourra trouver un
prolongement $g$ de~$g_k$ \`a $Y$. La condition de nullit\'e des
$\H^1$ est v\'erifi\'ee en particulier si $Y_0$ est affine. On
trouve donc:

\begin{corollaire}
\label{III.5.5}
Dans
l'\'enonc\'e du th\'eor\`eme~\ref{III.3.1}, on obtient
une condition n\'ecessaire et suffisante \'equivalente aux autres
en supposant que le~$Y^\prime$ qui intervient dans \textup{(ii)} (ou \textup{(ii~bis)})
est affine, et en exigeant l'existence d'un \emph{prolongement global}
$g$ de~$g_0$ \`a tout~$Y^\prime$.
\end{corollaire}

On notera que la d\'emonstration de~\ref{III.3.1} ne pouvait pas
nous donner ce r\'esultat directement.

Un cas important est celui o\`u $Y$ est \emph{plat} sur~$S$, alors
on a donc
\begin{equation*}
\gr^n(\cal{O}_Y)=\gr^n(\cal{O}_S)\otimes_{\cal{O}_{S_0}}\cal{O}_{Y_0}
\end{equation*}
et lorsque de plus les $\gr^n(\cal{O}_S)$ sont \emph{localement libres
sur}~$S$, on trouve
\begin{equation*}
\cal{G}_n=
\SheafHom_{\cal{O}_{Y_0}}(g_0^*(\it{\Omega}_{X_0/S_0}^1),\cal{O}_{Y_0})
\otimes_{\cal{O}_{S_0}}\gr^n(\cal{O}_S),
\end{equation*}
ou encore, si $\it{\Omega}_{X_0/S_0}^1$ est lui aussi localement libre
(par exemple $X$ lisse sur~$S$)
\begin{equation*}
\cal{G}_n=g_0^*(\goth{g}_{X_0/S_0})
\otimes_{\cal{O}_{S_0}}\gr^{n}(\cal{O}_S)
\end{equation*}
Si par exemple $S$ est affine d'anneau affine~$A$, $\cal{I}$ \'etant
d\'efini par un id\'eal $I$ de~$A$, on trouve
\begin{equation*}
\H^i(Y_0,\cal{G}_n)=\H^i(Y_0,\cal{G}_0) \otimes_A\gr_I^n(A)
\end{equation*}
pour tout $i$ (en effet, la question est locale sur~$S_0$, et on est
ramen\'e au cas o\`u on tensorise par un Module libre).
\emph{Dans ce cas, la nullité de~$\H^1(Y_0,\cal{G}_0)$ implique
que toutes les obstructions aux prolongements successifs de~$g_n$ sont
nulles}. On obtient donc:

\begin{corollaire}
\label{III.5.6}
Soient $(S,X,Y,\cal{I},g_n)$ comme plus haut, supposons de plus $X$
lisse sur~$S$ et $Y$ plat sur~$S$, enfin $S$ affine, et les
$\gr^n(\cal{O}_S)=\cal{I}^n/\cal{I}^{n+1}$ localement libres. Alors
l'obstruction \`a construire~$g_{n+1}$ se trouve dans
$\H^1(Y_0,\cal{G}_0)\otimes_A\gr_I^{n+1}(A)$ (o\`u~$A$ est l'anneau
de~$S$, $I$ l'id\'eal de~$A$ d\'efinissant $\cal{I}$), en posant
\begin{equation*}
\cal{G}_0=g_0^*(\goth{g}_{X_0/S_0})
\end{equation*}
Si $\H^1(Y_0,\cal{G}_0)=0$, alors $g_n$ peut se prolonger en un
$\widehat{S}$-morphisme $\widehat{g}\colon\widehat{Y}\to\widehat{X}$.
\end{corollaire}

Bien entendu, ce r\'esultat resterait valable tel quel, si on
partait, au lieu de~$S$-pr\'esch\'emas ordinaires $X$ et~$Y$, de
$\widehat{S}$-pr\'esch\'emas formels $\widehat{\cal{I}}$-adiques
$\goth{X}$ et~$\goth{Y}$. Il permet de prouver par exemple que
certains sch\'emas formels propres sur un anneau local complet (par
exemple) sont en fait alg\'ebriques. En effet, proc\'edant comme
dans le lemme~\ref{III.4.2}, on trouve:

\begin{corollaire}
\label{III.5.7} Sous les conditions de~\ref{III.5.6}, si $g_0$ est
un isomorphisme il en est de m\^eme de~$\widehat{g}$.
\end{corollaire}

\noindent
(N.B. le m\^eme r\'esultat vaut pour les immersions ferm\'ees).

On obtient ainsi:

\begin{proposition}
\label{III.5.8} Soient $A$ un anneau local complet d'id\'eal
maximal $\goth{m}$,
de
corps r\'esiduel~$k$, soient $\goth{X}$ et $\goth{Y}$ deux
pr\'esch\'emas formels $\goth{m}$-adiques sur~$A$, plats sur~$A$
(\ie pour tout~$n$, $X_n$ et $Y_n$ sont plats
sur~$A_n=A/\goth{m}^{n+1}$), on suppose $X_0=\goth{X}\otimes_Ak$ lisse
sur~$k$ et $\H^1(X_0,\goth{g}_{X_0/k})=0$. Alors tout $k$-isomorphisme
de~$Y_0$ sur~$X_0$ se prolonge en un $A$-isomorphisme de~$\goth{Y}$
sur~$\goth{X}$; ce prolongement est unique si de plus
$\H^0(X_0,\goth{g}_{X_0/k})=0$.
\end{proposition}

Cela donne en particulier un r\'esultat d'\emph{unicité} de
pr\'esch\'ema formel lisse sur~$A$ se r\'eduisant suivant un
pr\'esch\'ema $X_0$ donn\'e (moyennant
$\H^1(X_0,\goth{g}_{X_0/k})=0$). De plus, si~$\goth{X}$ et~$\goth{Y}$
proviennent de sch\'emas ordinaires propres sur~$A$, soient $X$
et~$Y$, alors on sait d'apr\`es le th\'eor\`eme d'existence de
faisceaux en G\'eom\'etrie formelle (\cf expos\'e au
S\'eminaire Bourbaki, \No 182\footnote{\Cf EGA III 5.4.1 pour la
d\'emonstration.}) qu'il y a correspondance biunivoque entre les
$A$-isomorphismes $Y\isomto X$ et le $A$-isomorphismes des
compl\'et\'es formels, donc

\begin{corollaire}
\label{III.5.9} L'\'enonc\'e pr\'ec\'edent~\ref{III.5.8} reste
valable en y rempla\c cant $\goth{X}$ et $\goth{Y}$ par des
$A$-sch\'emas ordinaires $X$ et~$Y$, \emph{propres} sur~$A$.
\end{corollaire}

Enfin, lorsque $\goth{X}$ est un sch\'ema formel propre sur~$A$, et
que $\goth{Y}$ est de la forme~$\widehat{Y}$ o\`u~$Y$ est un
sch\'ema ordinaire propre sur~$A$, alors la
proposition~\ref{III.5.8} donne des conditions suffisantes pour qu'on
puisse trouver un isomorphisme de~$\goth{X}$ avec~$\widehat{Y}$, donc
pour que le sch\'ema formel $\goth{X}$ soit en fait
\og alg\'ebrique\fg (\ie isomorphe \`a un~$\widehat{X}$, $X$ un
sch\'ema ordinaire propre sur~$A$, lequel sera alors canoniquement
d\'etermin\'e). C'est ce qui a lieu notamment si $X_0=\PP_k^r$ (ou
plus g\'en\'eralement, si $X_0$ est un sch\'ema de
Severi-Brauer,
\ie devient
isomorphe \`a l'espace projectif type sur la cl\^oture
alg\'ebrique de~$k$): tout sch\'ema formel propre et plat sur~$A$,
de fibre~$\PP_k^r$, est alg\'ebrisable, et de fa\c con plus
pr\'ecise est isomorphe au compl\'et\'e formel $\goth{m}$-adique
de~$\PP_A^r$. En particulier (gr\^ace au \og th\'eor\`eme
d'existence\fg) tout sch\'ema ordinaire propre sur~$A$, de
fibre~$\PP_k^r$, est isomorphe \`a~$\PP_A^r$ ($A$ \'etant un
anneau local complet). Utilisant la th\'eorie de la descente, on
peut prouver que si~$A$ n'est pas complet, $X$ devient isomorphe
\`a~$\PP^r$ en faisant une extension $A\to A^\prime$ finie et
\'etale de la base (et sous cette forme, le r\'esultat reste
valable pour une fibre qui est un sch\'ema de
Severi-Brauer).

\section{Prolongement infinit\'esimal global des $S$-sch\'emas lisses}
\label{III.6}

Sous les conditions du th\'eor\`eme~\ref{III.4.1}, on se propose
de chercher s'il existe un pr\'esch\'ema~$X$ lisse sur~$Y$ tel que
$X\times_YY_0$ soit $Y_0$-isomorphe \`a~$X_0$, sachant qu'un tel
sch\'ema \og existe localement sur~$X_0$\fg. Reprenant la m\'ethode
de construction de proche en proche, on est conduit \`a
remplacer~$Y$ par la lettre~$S$, \`a supposer qu'on se donne un
sous-pr\'esch\'ema ferm\'e $S_0$ de~$S$ d\'efini par un
faisceau d'id\'eaux~$\cal{I}$, (qu'il n'est plus n\'ecessaire de
supposer localement nilpotent), \`a introduire les
sous-pr\'esch\'emas ferm\'es~$S_n$ de~$S$ d\'efinis par
les~$\cal{I}^{n+1}$, et \`a supposer qu'on s'est donn\'e un
sous-pr\'esch\'ema~$X_{n}$ lisse sur~$S_n$. On se propose de
trouver un $S_{n+1}$-pr\'esch\'ema~$X_{n+1}$ \og qui se r\'eduit
suivant~$X_{n}$\fg, \ie muni d'un isomorphisme
\begin{equation*}
X_{n+1}\times_{S_{n+1}}S_n\isomfrom X_n
\end{equation*}
qui soit \emph{lisse} sur~$S_{n+1}$ (ou, ce qui revient au m\^eme
par II~\ref{II.2.1}, \emph{plat} sur~$S_{n+1}$). Comme nous l'avons
signal\'e dans \No \ref{III.4}, une telle donn\'ee revient \`a la
donn\'ee d'un faisceau d'alg\`ebres~$\cal{B}$ sur~$f^{-1}(\cal{O}_{S_{n+1}})$ (o\`u $f$ est l'application continue
sous-jacente au morphisme structural $X_n\to S_n$), muni d'une
augmentation $\cal{B}\to\cal{O}_{X_n}$ compatible avec l'augmentation
$f^{-1}(\cal{O}_{S_{n+1}})\to f^{-1}(\cal{O}_{S_n})$, et satisfaisant
\`a deux conditions (a) et (b) que nous ne r\'e\'ecrirons pas,
nous bornant \`a noter qu'elles sont de \emph{nature locale} sur
l'espace topologique sous-jacent \`a~$X_n$. On sait
d'apr\`es~\ref{III.4.1} qu'une solution existe localement. Elle est
de plus unique \`a isomorphisme (non unique) pr\`es, du moins
localement. Commen\c cons par pr\'eciser ce point:

\begin{proposition}
\label{III.6.1}
Soit
$X_{n+1}$ sur~$S_{n+1}$ se r\'eduisant suivant $X_n$
sur~$S_{n}$. Alors le faisceau (sur l'espace topologique sous-jacent
\`a~$X_{n}$, ou encore \`a~$X_0$) des $S_{n+1}$-automorphismes
de~$X_{n+1}$ qui induisent l'identit\'e sur~$X_n$ est canoniquement
isomorphe~\`a
\begin{equation*}
\cal{G}=\goth{g}_{X_0/S_0}\otimes_{\cal{O}_{S_0}}
\gr_{\cal{I}}^{n+1}(\cal{O}_S)
\end{equation*}
(en tant que faisceau de groupes).
\end{proposition}

En effet, par~\ref{III.5.4} et~\ref{III.4.2} ce faisceau est un
faisceau principal homog\`ene sous~$\cal{G}$. Comme il est muni
d'une section privil\'egi\'ee (l'automorphisme identique
de~$X_{n+1}$), il s'identifie donc comme faisceau d'ensembles
\`a~$\cal{G}$. Il faut v\'erifier que cette identification est
compatible avec les structures de groupe. C'est facile, et d'ailleurs
un cas particulier d'un r\'esultat plus g\'en\'eral sur la
compatibilit\'e des structures de fibr\'es principaux,
dans~\ref{III.5.1} et~\ref{III.5.3}, avec la composition des
morphismes (r\'esultat que nous n'\'enon\c cons pas ici, mais
qui se doit de figurer dans le hyperplodoque).

En particulier, le faisceau sur~$X_0$ des germes d'automorphismes
de~$X_{n+1}$ (avec les structures explicit\'ees) est
\emph{commutatif}. Il s'ensuit que si~$X_{n+1}^\prime$ est une autre
solution du probl\`eme, isomorphe \`a~$X_{n+1}$ au-dessus de
l'ouvert~$U$ de~$X_0$, alors l'isomorphisme de~$\SheafAut(X_{n+1})|U$
sur~$\SheafAut(X_{n+1}^\prime)|U$ d\'eduit par transport de
structure d'un isomorphisme $X_{n+1}|U\isomto X_{n+1}^\prime|U$,
\emph{ne dépend pas} du choix de ce dernier. (Ce n'est d'ailleurs
autre que l'isomorphisme identique de~$\cal{G}$, lorsque on identifie
l'un et l'autre faisceau d'automorphismes \`a~$\cal{G}$ gr\^ace
\`a~\ref{III.6.1}).

On d\'eduit de~\ref{III.6.1}:

\begin{corollaire}
\label{III.6.2}
Soient~$X_{n+1}$, $X_{n+1}'$ lisses sur~$S_{n+1}$ et \og se
r\'eduisant suivant~$X_n$\fg. Alors le faisceau (sur l'espace
sous-jacent \`a~$X_0$) des~$S_{n+1}$-isomorphismes de~$X_{n+1}$
sur~$X_{n+1}'$ induisant l'identit\'e sur~$X_n$, est de fa\c con
naturelle un faisceau principal homog\`ene sous~$\cal{G}$.
\end{corollaire}

Cela exprime en effet que~$X_{n+1}$ et~$X_{n+1}'$ sont isomorphes
localement, et que le faisceau des germes d'automorphismes du premier
est~$\cal{G}$.

Notons maintenant qu'en vertu de~\ref{III.4.1}, on peut toujours
trouver un recouvrement~$(U_i)$
de~$X_n$ par des ouverts (qu'on peut supposer affines), et pour
tout~$i$ un sch\'ema lisse~$X^i$ sur~$S_{n+1}$ se r\'eduisant
suivant~$U_i$.
Supposons pour simplifier~$X_n$ \emph{séparé}, donc
les~$U_{ij}=U_i\cap U_j$ sont encore des ouverts \emph{affines}
de~$X_n$. Comme le~$\H^1$ d'un tel ouvert, \`a valeurs dans le
faisceau quasi-coh\'erent~$\cal{G}$, est nul, on en d\'eduit par
le
corollaire~\ref{III.6.2} que~$X^i|U_{ij}$ est isomorphe
\`a~$X^j|U_{ij}$; soit
$$
f_{ji}:X^i|U_{ij}\isomto X^j|U_{ij}
$$
un tel isomorphisme. Il est d\'etermin\'e \`a une section
pr\`es de~$\cal{G}$ sur~$U_{ij}$. Posons, pour tout
triplet
d'indices:
$$
f_{ji}^{(k)}=f_{ji}|U_{ijk}\quad\text{où }U_{ijk}=U_i\cap U_j\cap
U_k.
$$
Si on avait
\begin{equation*}
\label{eq:III.6.1}
\tag{1} {f_{kj}^{(i)}f_{ji}^{(k)}=f_{ki}^{(j)},}
\end{equation*}
il s'ensuivrait que les~$X^i$ se \og recollent\fg par les~$f_{ji}$, donc
qu'ils d\'efinissent une solution~$X=X_{n+1}$ du probl\`eme
cherch\'e. Une telle solution existe plus g\'en\'eralement, si
on peut modifier les~$f_{ji}$ en des~$f_{ji}'$:
\begin{equation*}
\label{eq:III.6.2}
\tag{2} {f_{ji}'=f_{ji}g_{ji}\quad (g_{ji}\in\Gamma(U_{ij},\cal{G}))}
\end{equation*}
de telle fa\c con que les~$f_{ji}'$ satisfassent la condition de
transitivit\'e ci-dessus. Cette condition suffisante pour
l'existence d'une solution est aussi n\'ecessaire, comme on voit en
se rappelant qu'une telle solution~$X$ doit, sur chaque~$U_i$,
\^etre isomorphe \`a~$X^i$, ce qui permet donc de choisir des
isomorphismes
$$
f_i:X|U_i\isomto X^i
$$
et de d\'efinir des
$$
f_{ji}'=(f_j|U_{ij})(f_i|U_{ij})^{-1}:X^i|U_{ij}\isomto X^j|U_{ij}
$$
satisfaisant la condition de recollement.

Or posons
\begin{equation*}
\label{eq:III.6.3}
\tag{3}
{f_{ijk}=(f_{ki}^{(j)})^{-1}f_{kj}^{(i)}f_{ji}^{(k)},}
\end{equation*}
c'est un automorphisme de~$X^i|U_{ijk}$, que nous identifierons \`a
une section de~$\cal{G}$ gr\^ace \`a~\ref{III.6.1}. On constate
que c'est un $2$-\emph{cocycle} $f$ du recouvrement ouvert
$\cal{U}=(U_i)$, \`a coefficients dans~$\cal{G}$, par un petit
calcul formel laiss\'e au lecteur. Le m\^eme calcul montre que
moyennant~\eqref{eq:III.6.2}, la condition de
recollement~\eqref{eq:III.6.1} \emph{pour les} $f_{ij}'$ \'equivaut
\`a la formule
\begin{equation*}
\label{eq:III.6.4}
\tag{4} {f=dg,}
\end{equation*}
o\`u $g=(g_{ij})$ est consid\'er\'e comme une $1$-cocha\^ine
de~$\cal{U}$ \`a coefficients dans~$\cal{G}$. Donc \emph{la
condition nécessaire et suffisante pour l'existence d'une solution
du problème est que la classe de cohomologie dans
$\H^2(\cal{U},\cal{G})$ définie par le cocycle~\eqref{eq:III.6.3}}
soit nulle. Notons d'ailleurs que puisque~$\cal{U}=(U_i)$ est un
recouvrement affine de~$X_0$ qui est un \emph{schéma},
$\H^2(\cal{U},\cal{G})$ s'identifie \`a $\H^2(X_0,\cal{G})$. Il est
imm\'ediat d'ailleurs que la classe de cohomologie ainsi obtenue
dans $\H^2(X_0,\cal{G})$ ne d\'epend pas du recouvrement affine
consid\'er\'e. On l'appellera la \emph{classe d'obstruction au
prolongement de~$X_n$ en un schéma~$X_{n+1}$ lisse sur~$S_{n+1}$}.

Supposons cette obstruction nulle. Alors un raisonnement esquiss\'e
plus haut montre que toute solution~$X=X_{n+1}$ est isomorphe \`a
une solution obtenue par recollement \`a partir
d'isomorphismes~$f_{ji}'$, qu'on peut supposer sous la
forme~\eqref{eq:III.6.2}, la condition de recollement n'\'etant
autre que~\eqref{eq:III.6.3}. L'ensemble des $g$ admissibles est donc
un espace principal homog\`ene sous le groupe
$\mathrm{Z}^1(\cal{U},\cal{G})$ des $1$-cocycles de~$\cal{U}$ \`a
coefficients dans~$\cal{G}$. De plus, on constate tout de suite que
\emph{deux cochaînes $g$ et $g'$} (telles que $dg=dg'=f$)
\emph{définissent des solutions isomorphes si et seulement si le
cocycle $g-g'$ est de la forme $dh$}, o\`u
$h=(h_i)\in\mathrm{C}^0(\cal{U},\cal{G})$. On trouve donc:

\begin{theoreme}
\label{III.6.3}
\index{obstruction au prolongement d'un sch\'ema en un sch\'ema lisse|hyperpage}Soient $(S,\cal{I},X_n)$ comme ci-dessus, $X_n$ \'etant suppos\'e
s\'epar\'e\footnote{Cette condition est en fait inutile, et on
peut \'eviter les calculs de cocycles plus haut. \Cf le livre de J.
\textsc{Giraud}, Cohomologie Non Ab\'elienne (\`a para\^itre dans
Springer Verlag 1971). Comparer remarques~\ref{III.6.4}.}. Alors on
peut d\'efinir canoniquement une classe d'obstruction dans
$\H^2(X_0,\cal{G})$ (o\`u~$\cal{G}$ est d\'efini
dans~\ref{III.6.1}), dont l'annulation est n\'ecessaire et
suffisante pour l'existence d'un sch\'ema~$X_{n+1}$ lisse
sur~$S_{n+1}$ se r\'eduisant suivant~$X_n$. Si cette obstruction
est nulle, alors l'ensemble des classes, \`a isomorphisme pr\`es
(induisant l'identit\'e sur~$X_n$)
de~$S_{n+1}$-pr\'esch\'emas~$X_{n+1}$ se r\'eduisant
suivant~$X_n$, est de fa\c con naturelle un espace principal
homog\`ene sous $\H^1(X_0,\cal{G})$.
\end{theoreme}

\begin{remarques}
\label{III.6.4}
\`A partir de~\ref{III.6.1}, les raisonnements faits ici sont purement
formels, et se transcrivent avantageusement dans le cadre des
cat\'egories locales,
ou m\^eme des cat\'egories fibr\'ees g\'en\'erales. La
classe d'obstruction \`a l'existence d'un objet \og global\fg d'une
cat\'egorie (dont on peut trouver un objet \og localement\fg, et dont
deux objets sont toujours \og isomorphes localement\fg, le groupe des
automorphismes de tout objet \'etant commutatif) ainsi obtenu dans
le contexte g\'en\'eral, contient comme cas particulier le
\og deuxi\`eme homomorphisme bord\fg dans une suite exacte de faisceaux
de groupes non n\'ecessairement commutatifs, (\'etudi\'e par
exemple par Grothendieck dans Kansas ou
T\^ohoku).
Le calcul b\'eb\^ete par cocycles fait ici doit donc \^etre regard\'e
comme un pis-aller, d\^u \`a la non existence d'un texte de
r\'ef\'erence satisfaisant.
\end{remarques}

\subsection{}
\label{III.6.5}
On notera que dans~\ref{III.6.3}, il n'y a pas en g\'en\'eral
d'\'el\'ement privil\'egi\'e dans l'espace principal
homog\`ene envisag\'e sous $\H^1(X_0,\cal{G})$. Cela se traduit
notamment par le fait que l'on obtient (localisant sur~$S$) un
faisceau principal homog\`ene sur~$S_0$, de groupe structural
$\R^1f_*(\cal{G})$, qui n'est pas n\'ecessairement trivial, \ie qui
d\'efinit une classe de cohomologie dans
$\H^1(S_0,\R^1f_*(\cal{G}))$ qui n'est pas n\'ecessairement nulle.
(Lorsque l'on suppose que la classe $d\in\H^2(X_0,\cal{G})$ n'est pas
nulle, mais nulle \og localement au-dessus de~$S$\fg, \ie d\'efinit
une section nulle de~$\R^2f_*(\cal{G})$, \ie un \'el\'ement nul
dans $\H^0(S_0,\R^2f_*(\cal{G}))$).

\subsection{}
\label{III.6.6}
On ne sait pour l'instant \`a peu pr\`es rien sur le m\'ecanisme
alg\'ebrique g\'en\'eral des classes de cohomologie introduites
dans ce num\'ero et ses relations avec le num\'ero
pr\'ec\'edent, et on ne sait rien en dire de pr\'ecis dans les
cas particuliers les plus simples, tel le cas des sch\'emas
ab\'eliens sur des anneaux artiniens\footnote{\label{III.6.6.p24}On sait maintenant que cette obstruction est toujours nulle dans ce
cas \lcrochetbf ajout\'e en 2003 (MR): \cf F.\, Oort, Finite group schemes, local moduli for abelian varieties and
lifting problems, Algebraic Geometry Oslo 1970, Wolters-Noordhoff, 1972, p.\,223--254\rcrochetbf.}. On esp\`ere qu'il se trouvera des gens pour chiader la
question, qui semble particuli\`erement int\'eressante. Elle est
intimement li\'ee en particulier \`a la \og th\'eorie des
modules\fg des structures alg\'ebriques.

\begin{corollaire}
\label{III.6.7}
Supposons que $\H^2(X_0,\cal{G})=0$, alors un~$X_{n+1}$ existe, et il
est unique \`a isomorphisme pr\`es si de plus
$\H^1(X_0,\cal{G})=0$.
\end{corollaire}

En particulier, on en conclut, en proc\'edant de proche en proche
(et remarquant qu'un sch\'ema affine est acyclique pour un faisceau
quasi-coh\'erent):

\begin{corollaire}
\label{III.6.8}
Sous
les conditions du th\'eor\`eme~\ref{III.4.1}, si~$X_0$ est
affine, alors il existe un~$X$ lisse sur~$Y$ se r\'eduisant
suivant~$X_0$, et cet~$X$ est unique \`a isomorphisme (non unique)
pr\`es.
\end{corollaire}

On notera que la d\'emonstration directe du th.\, \ref{III.4.1} ne
pouvait nous donner ce r\'esultat.

\begin{corollaire}
\label{III.6.9}
Sous les conditions de~\ref{III.6.3}, supposons~$S$ affine
d'anneau~$A$, $\cal{I}$ d\'efini par un id\'eal~$I$ de~$A$, enfin
les $\gr^n_{\cal{I}}(\cal{O}_S)=\cal{I}^n/\cal{I}^{n+1}$ localement
libres. Alors $\H^i(X_0,\cal{G})$ s'identifie \`a
$\H^i(X_0,\cal{G}_0)\otimes_A\gr^{n+1}_I(A)$, o\`u
$$
\cal{G}_0=\goth{g}_{X_0/S_0},
$$
donc la classe d'obstruction au prolongement de~$X_n$ se trouve dans
$\H^2(X_0,\cal{G}_0)\otimes_A\gr^{n+1}_I(A)$, et si elle est nulle,
l'ensemble des classes (\`a isomorphisme pr\`es) de solutions est
un espace principal homog\`ene sous
$\H^1(X_0,\cal{G}_0)\otimes_A\gr^{n+1}_I(A)$.
\end{corollaire}

En particulier:

\begin{corollaire}
\label{III.6.10}
Sous les conditions de~\ref{III.6.9} supposons
$$
\H^2(X_0,\goth{g}_{X_0/S_0})=0,
$$
alors il existe un sch\'ema formel $\hat{I}$-adique
$\goth{X}$
sur le compl\'et\'e formel $\cal{I}$-adique~$\hat{S}$ de~$S$, qui
soit \og lisse sur~$S$\fg (\ie les~$\goth{X}_p$ sont lisses sur
les~$S_p$) et qui se r\'eduise suivant~$X_n$, \ie muni d'un
isomorphisme
$$
\goth{X}\times_S S_n\isomfrom X_n.
$$
Si de plus $\H^1(X_0,\goth{g}_{X_0/S_0})=0$, alors un tel~$\goth{X}$
est unique \`a isomorphisme pr\`es.
\end{corollaire}

En effet, on construit de proche en proche $X_{n+1}$,~$X_{n+2}$, etc.,
d'o\`u~$\goth{X}$ en passant \`a la limite inductive des~$X_i$.
L'assertion d'unicit\'e figure d\'ej\`a au \no
pr\'ec\'edent.

\section[Application \`a la construction de sch\'emas formels...]{Application \`a la construction de sch\'emas formels
et de sch\'emas ordinaires lisses sur un anneau local complet~$A$}
\label{III.7}
Les r\'esultats du \no pr\'ec\'edent permettent parfois de
prouver l'existence
d'un sch\'ema formel $\goth{m}$-adique sur un tel anneau, se
r\'eduisant suivant un sch\'ema lisse~$X_0$ sur~$k$
donn\'e. Distinguons deux cas:

a) \emph{$A$ est \og d'égales caractéristiques\fg} (c'est le cas
en particulier si~$k$ est de caract\'eristique~$0$). Alors on sait
qu'il existe un \emph{sous-corps de représentants de}~$A$, \ie
un sous-corps~$k'$ tel que~$A\to k$ induise un isomorphisme~$k'\isomto
k$. \emph{Alors il existe même un schéma ordinaire lisse
sur~$A$ se réduisant suivant~$X_0$}, savoir~$X=X_0\otimes_kA$, $A$~\'etant consid\'er\'e comme une alg\`ebre sur~$k$ gr\^ace
\`a l'homomorphisme~$k\to k'\to A$ d\'efini par~$k'$. Il faut
cependant noter que cette construction n'est pas \og naturelle\fg; il est
facile de se convaincre (d\'ej\`a dans le cas
o\`u~$A=k[t]/(t^2)$, alg\`ebre des nombres duaux) qu'un autre
homomorphisme de rel\`evement~$k\to A$ (d\'efini en l'occurrence
par une d\'erivation absolue de~$k$ dans lui-m\^eme) d\'efinit
un~$X'$ sur~$A$ qui en g\'en\'eral \emph{n'est pas isomorphe
à~$X$} (si~$\H^1(X_0,\goth{g}_{X_0/k})\neq0$). Il serait
d'ailleurs int\'eressant d'\'etudier (pour~$k$ de
caract\'eristique~$0$, ou non parfait et de
caract\'eristique~$p>0$) quels sont
les~$X$ lisses sur~$A$ qu'on obtient ainsi, \`a quelle condition
deux homomorphismes~$k\to A$ d\'efinissent des $A$-sch\'emas
isomorphes. N\'eanmoins, l'existence de~$k'$ suffit \`a
entra\^iner que la premi\`ere obstruction au rel\`evement
de~$X_0$, qui est
dans~$\H^2(X_0,\goth{g}_{X_0/k})\otimes_k\goth{m}/\goth{m}^2$, est
n\'ecessairement nulle. Bien entendu, quand on a relev\'e
alors~$X_0$ en~$X_1$ lisse sur~$A/\goth{m}^2$, la nouvelle obstruction
\`a la construction de~$X_2$ ne sera en g\'en\'eral pas nulle:
elle sera fonction d'un \'el\'ement variable dans un certain
espace principal homog\`ene
sous~$\H^1(X_0,\cal{G}_0)\otimes\goth{m}/\goth{m}^2$ et se trouve
dans~$\H^2(X_0,\cal{G}_0)\otimes\goth{m}^2/\goth{m}^3$: il
conviendrait d'\'etudier la situation de fa\c con
d\'etaill\'ee\footnote{Elle est sans doute d\'ecrite par
l'op\'eration crochet de Kodaira-Spencer (\cf S\'eminaire Cartan,
1960/61, Exp.\, 4).}.

b) \emph{$A$ est d'inégales caractéristiques.} Dans ce cas, on
ignore tout, (sauf si par chance $\H^2(X_0,\goth{g}_{X_0/k})=0$,
auquel cas on peut construire un sch\'ema formel $\goth{m}$-adique
simple sur~$A$ se r\'eduisant suivant~$k$). M\^eme
si~$A=\ZZ/p^2\ZZ$ et si~$X_0$ est un sch\'ema \og ab\'elien\fg de
dimension~$2$, on ne sait pas si on peut le relever en un~$X=X_1$ lisse
sur~$A$\footnote{C'est maintenant prouv\'e, \cf \
note~\ref{III.6.6.p24} page~\pageref{III.6.6.p24}.}, d'autre part, on
n'a pas d'exemple d'un~$X_0$ dont on ait prouv\'e qu'il ne provient
pas d'un sch\'ema ordinaire~$X$ lisse sur~$A$. (J'ai l'impression
que cela doit exister, avec~$X_0$ une surface
projective)\footnote{\label{nIII.3}Un tel exemple a \'et\'e depuis
construit par
J-P.\, Serre (Proc.\ Nat.\ Acad.\ Sc.\ USA  \textbf{47} (1961),
\No 1, p.\,108--109), du moins pour certaines
dimensions. D.\, Mumford a donn\'e un exemple (non publi\'e) avec
une \emph{surface} alg\'ebrique.}. Signalons simplement que
d'apr\`es le th\'eor\`eme de Cohen, il existe un $p$-anneau de
Cohen~$B$ de corps r\'esiduel~$k$ et un homomorphisme~$B\to A$
induisant l'isomorphisme identique sur les corps r\'esiduels; par
suite, le r\'esultat \og le plus fort\fg de rel\`evement serait
obtenu en prenant pour~$A$
un $p$-anneau de Cohen: s'il existe une solution (ordinaire ou
formelle) au-dessus d'un tel anneau, il en existe une au-dessus de
tout anneau local complet de corps r\'esiduel~$k$. En particulier,
comme pour un $p$-anneau de Cohen~$\goth{m}/\goth{m}^2$ s'identifie
canoniquement \`a~$k$, on voit que \emph{pour tout schéma
lisse~$X_0$ sur un corps~$k$ de caractéristique~$p>0$, il existe
une classe de cohomologie dans~$\H^2(X_0,\goth{g}_{X_0/k})$}
premi\`ere obstruction au rel\`evement de~$X_0$ en un sch\'ema
lisse sur un $p$-anneau de Cohen; on ignore si elle peut \^etre non
nulle\footnote{Elle peut \^etre non nulle, comme signal\'e en
note~\ref{nIII.3} \`a la page~\pageref{nIII.3}.}.

M\^eme si on arrive de proche en proche \`a construire les~$X_n$
se r\'eduisant suivant~$X_0$, cela ne donne en g\'en\'eral qu'un
sch\'ema \emph{formel}~$\goth{X}$ lisse sur~$A$, se r\'eduisant
suivant~$X_0$. Lorsque~$X_0$ est propre sur~$A$, il reste la question
si~$\goth{X}$ est en fait alg\'ebrisable, pour pouvoir obtenir un
sch\'ema \emph{ordinaire} propre sur~$A$ et simple sur~$A$, se
r\'eduisant suivant~$X_0$. Le seul crit\`ere connu (signal\'e
dans le S\'eminaire Bourbaki, et qui figure dans les \'El\'ements,
Chap.\, III, 4.7.1) est le suivant: si~$\goth{X}$ est propre sur~$A$, et
si~$\cal{L}$ est un faisceau inversible sur~$\goth{X}$ tel que le
faisceau induit~$\cal{L}_0$ sur~$X_0$ soit ample (\ie une puissance
tensorielle convenable~$\cal{L}_0^{\otimes n}$, $n>0$, provient d'une
immersion projective de~$X_0$) alors il existe un sch\'ema~$X$
projectif sur~$A$, et un faisceau inversible ample sur~$X$, tels
que~$(\goth{X},\cal{L})$ s'en d\'eduise par compl\'etion
$\goth{m}$-adique. Cela nous am\`ene donc, \'etant donn\'e un
faisceau localement libre~$\cal{E}_0$ sur~$X_0$ (que nous choisirons
inversible ample pour notre propos), de le prolonger en un faisceau
localement libre~$\cal{E}$ sur~$\goth{X}$. Pour ceci, on est
ramen\'e \`a construire de proche en proche des faisceaux
localement libres~$\cal{E}_n$ sur les~$X_n$. La discussion est toute
analogue \`a celle du \No \ref{III.6}, (\cf remarque~\ref{III.6.4}), le r\^ole
essentiel \'etant jou\'e par le \emph{faisceau des automorphismes}
d'un~$\cal{E}_{n+1}$ qui induisent l'identit\'e sur~$\cal{E}_n$: on
montre aussit\^ot que ce faisceau s'identifie \`a
\begin{equation*}
\label{eq:III.7.*}\tag{$*$}{}
\cal{G}= \SheafHom_{\cal{O}_{X_0}}(\cal{E}_0,
\cal{E}_0\otimes\gr_{\cal{I}}^{n+1}(\cal{O}_X))=
\SheafHom_{\cal{O}_{X_0}}(\cal{E}_0,
\cal{E}_0)\otimes\gr_{\cal{I}}^{n+1}(\cal{O}_X)
\end{equation*}
qui est encore un faisceau en groupes commutatifs. On trouve:

\begin{proposition}\label{III.7.1}
Soit~$S$ un pr\'esch\'ema muni d'un faisceau quasi-coh\'erent
d'id\'eaux~$\cal{I}$, $X$ un pr\'esch\'ema sur~$S$, $S_n$ le
sous-pr\'esch\'ema de~$S$ d\'efini par~$\cal{I}^{n+1}$,
et~$X_n=X\times_SS_n$ (pour tout entier~$n$). Soit~$\cal{E}_n$ un
faisceau localement libre
sur~$X_n$, on se propose de le prolonger en un faisceau localement
libre~$\cal{E}_{n+1}$ sur~$X_{n+1}$. Alors~$\cal{E}_n$ d\'efinit une
classe d'obstruction canonique dans~$\H^2(X_0,\cal{G})$,
o\`u~$\cal{G}$ est le faisceau quasi-coh\'erent donn\'e par la
formule ci-dessus, classe dont l'annulation est n\'ecessaire et
suffisante pour l'existence d'un~$\cal{E}_{n+1}$
prolongeant~$\cal{E}_n$. Si cette classe est nulle, alors l'ensemble
des classes, \`a isomorphisme pr\`es (induisant l'identit\'e
sur~$\cal{E}_n$) de solutions~$\cal{E}_{n+1}$, est un espace principal
homog\`ene sous~$\H^1(X_0,\cal{G})$.
\end{proposition}

Cette proposition donne lieu aux corollaires habituels. Signalons
seulement que si~$X$ est \emph{plat} sur~$S$, alors on peut \'ecrire
$$
\cal{G}= \SheafHom_{\cal{O}_{X_0}}(\cal{E}_0, \cal{E}_0)
\otimes_{\cal{O}_{S_0}} \gr_{\cal{I}}^{n+1}(\cal{O}_S)
$$
d'o\`u, si~$S$ est affine d'anneau~$A$ et si
les~$\cal{I}^n/\cal{I}^{n+1}$ sont localement libres, la condition
suffisante
$$
\H^2(X_0,\cal{G}_0)=0 \qquad\text{avec}\qquad
\cal{G}_0=\SheafHom_{\cal{O}_{X_0}}(\cal{E}_0, \cal{E}_0)
$$
pour l'existence d'un~$\cal{E}_{n+1}$, donc de proche en proche pour
l'existence de prolongements successifs~$\cal{E}_m$ ($m=n,n+1,$
etc.).

Revenant \`a la situation de d\'epart, on trouve donc:

\begin{proposition}\label{III.7.2}
Soient~$A$ un anneau local complet, $\goth{X}$ un sch\'ema formel
propre et plat sur~$A$, tel que~$X_0$ soit projectif et
que~$\H^2(X_0,\cal{O}_{X_0})=0$. Alors il existe un sch\'ema~$X$
projectif sur~$A$ dont le compl\'et\'e formel $\goth{m}$-adique
est isomorphe \`a~$\goth{X}$.
\end{proposition}

Conjuguant avec~\ref{III.6.10}, on trouve:

\begin{theoreme}\label{III.7.3}
Soient~$A$ un anneau local complet de corps r\'esiduel~$k$, $X_0$ un
sch\'ema projectif et lisse sur~$k$, tel que
$$
\H^2(X_0,\goth{g}_{X_0/k})=\H^2(X_0,\cal{O}_{X_0})=0
$$
Alors il existe un sch\'ema lisse et projectif~$X$ sur~$A$, se
r\'eduisant suivant~$X_0$.
\end{theoreme}

Plus g\'en\'eralement, si on se donne un~$X_n$ lisse
sur~$A_n=A/\goth{m}^{n+1}$ se r\'eduisant suivant~$X_0$, alors il
existe un~$X$ lisse et propre sur~$A$ et un
isomorphisme~$X\otimes_AA_n=X_n$.

\begin{corollaire}\label{III.7.4}
Toute courbe lisse et propre sur~$k$ provient par r\'eduction d'une
courbe lisse et propre sur~$A$.
\end{corollaire}

C'est ce r\'esultat qui sera l'outil essentiel (avec le
th\'eor\`eme d'existence de faisceaux en G\'eom\'etrie
Formelle) pour \'etudier le groupe fondamental de~$X_0$ par voie
transcendante.

